%! Author = lili
%! Date = 6/24/26

\untit{Inklusive Notizen Vorlesung 24. Juni 2026; Kapitel 16 in Lewin's Genes II}

Chaperone gibts auch im Mitochordium, aber das ist im Bild nicht drin.
Die liegen in der teilweise ungefalteten Form vor irgendwo und sind dann noch nicht aktiv.
Chaperone energetisieren auch den Transport

SRP würde Transitpeptit nicht erkennen, dafür ist Translocon da? Googeln

Wa sind antikörperund wo kommen die her und was für / wieso hat der verschiedene Spezifitäten sollte man schon wissen

``Histokompatibilitätskomplex'' ist auch Gewebeverträglichkeitskomplex

\defin{gewebeverträglichkeitskomplex}

\defin{innate_immunity}

``innate immunity'' ist das angeborene Immunität, die man nicht über das Immunsystem ``erlernen'' kann
Pflanzen und Insketen haben nur innate immunity

Abwehrreaktionen beruht auf bestimmten Erkennungsprozessen. Was erkannt werden kann ist bei innate immunity
vorgegeben, bei der erlernten Immunität nicht. Man nennt es auch angeborenen Resistenz, weil es ja an sich nicht der
gleiche Prozess ist.

\kfrage{innateimunity}

\defin{tlr}
tlr haben diese vorgegbenen Erkennungsfunktionen

\defin{mamp}
Si ehießen erst pamp, bis man festgestellt hat die sind nicht pathogen speziefisch, sondern mikrobenspezifisch.
reagieren auf Sachen, die ein Bakterium nicht so ganz losweredn kann wie bsp. Flagellen.

Flaggelen sind hochkonserviert undn haben hochkonservierte Proteine und die könen dann erkannt werden.

\defin{prr}

Elicitoren = Signalstoffe

% TODO auflisten was zu welchem Zellteil bzw. Organismus gehört

Folie 13 die Abbildung ist falsch, denn die Antikörper haben auf dem Bild nur eine Bindzngstelle, aber sie haben in
echt 2!
Nur deswegen können sie überhaupt als Antiköprer funktionieren, denn so können sie Antikörper verketten und
Komplexe bilden, die dann von \gls{macrophage}n gefressen werden können.
Die fressen nur Komplexe und keinen einzelnen Dingsies.

Folie 14 die ABb ist etwas besser, auch auch immer noch falsch.

%T-Zellen werden im Thymus differenziert. B-Zellen haben ihren Namen von V

T-Zellen sind für die zelluläre Immunität.

Bspw. wann kann ein Kaninchen immuniesieren und dessen Serum abnehmen und damit im Labor arbeiten, weil die
Antikörper da dann auch drin sind.

Gibt große allelische Diversität bei den irgendwas - Antigenen? Rezeptoren?

Heute kann man monoklonale Antikörper auch aus bakteriellen Bibliotheken machen, deshalb sind Mausexperimente nicht
mehr (unbedingt) nötig.


MAC / MHH Ausstattung ist das was bei DKMS Typisierungen geacht wird?

% 17

Das Immunsystem kann grundlegen gegen alles vorgehen, auch gegen sich selbst.
Damit das nicht passiert wid durch die klonale Deletion aussortiert, was das Selbst angreifen würde.

eine Zelle - eine Spezifität, damit das ganze Funktioniert mit der Selektion/Deletion

Abb F 18 oben das vorhandene Spektrum, dann wird eine ausgewählt, die nicht das Selbst angreift und die wird
vervielfältigt.

\frage{wieimmunsystemalle}

\frage{verschiedeneanikoerper}

Aufbau AK

2 Bindungsstellen
4 Polypeptide

Die Bindestelle ist nicht symmetrisch wie auf dem Bild (als V) sondern absolut identisch, auch stereospezifisch),
müsste eigentlich parallel.

\defin{leichte_kette}
\defin{schwere_kette}
\defin{vregion}
\defin{cregion}
\defin{hapten} kleines mölekül wa sman auf die eines größen koppelt um ein Antikörper für ein kleines zu machen.
Bspw. Zcker sind normalerweise nicht immunogen, aber wenn man das so macht, dann iwie schon.
\defin{epitope}
Bereich der gut erkannt wird auf etwas
\defin{antigene_determinante}
antigene determinante ist eine Art von Epitop?

% F 22
Diese Segmente sollte man nicht Gen nennen, sondern Segment, weil es nicht unserer Gendefinition entspricht. Es wird
halt erst zum Gen durch somatische Rekombination.

Keimbahn nur locus ohne Gene, nur segmente, dann som. rek. baut gene zusammen

\defin{lymphocyte_dna}

% TODO zusammebbau der leichten und schweren kette gegenüber stellen (VJ vs VDJ)

D wie diversity - mit drei komponenten kann man mehr kombinationen bilden als mit zweien

F 26 vH Gen ist auch wieder nur ein Segment, kein Gen

Wenn man sich wunder warum Abbildungen wie F 27 verschieden aussehen sollte man erstmal schaue, ob man sich Maus
oder Mensch anschaut.
Durchaus verschieden.

F 27 Pseudogene = Pseudogensegmente

Die Pseudogene tragen im Individuum nichts bei, haben keine Auswirkung

F 29 es gibt sowas wie in rechtes und linkes signal - huh?
Das was zeischen den beiden spacern liegt ist das, was verändert/herausgeschnitten wird  (bei F 31 ist das das was
zwischen den Dreiecken ist)

% NHEJ was ist das? Repariert iwie Doppelstrangbrüche?

F32 stellt den Prozess schöner da

Der Auswahlprozess, was in dem paired complex ausgewählt wird, ist sehr wichtig, essentiell, denn das wird ja dann
verwendet

% F 33
CTCF relevant

So eine schleife (rechts) kann nur entstehen, wenn die Bindestellen gegenläufig sind, deswegen ist in der mittleren
abb der erste linke nicht ´´hängen'' geblieben

% F 38
lg-lg- bekommt man praktisch nicht/kaum

Jede Zelle hat für jeden Loci sozusagen zwei Chancen, weil es bei jeder der beiden Kopien(?) klappen kann

% TODO Tabelle was welche kombi bedeutet bzw. ob es die gibt

\defin{hklasse}

F 44 Stichwort Immunisierungsschema - muss probiert werden, weil bisher nicht vorhersagbar

Plasmazellen entwickeln sich aus (Teilen?) der B-zellen

\defin{shm}

\frage{wokommtakher}

\defin{tcr}

\defin{mhc}
MHC ist das bei der Maus
\defin{hla}

F 50 das ist wohl eine Macrophage?
Anitgen-Präsenttaion nochmal googeln. Dieser Prozess passiert angeblich mit allem was die Phage aufbimmt oder was
irgendwas aufnimmt? Unklar.

F 51 abbildung links das wäre eigentlich ein größerer Komplex als nur zwei schwänzchen zu haben und auch mehr
entikörper mit mehr Antigenen würden ein Komplex bilden als das was da ist mit den 2:2

Auch normalerweise polyklonal und nicht nur einer


% Viren sind nicht im Serum? Warum?